\documentclass[11pt,letterpaper]{article}
\usepackage[margin=0.7in]{geometry}
\usepackage[T1]{fontenc}
\usepackage{mathptmx}
\usepackage[hidelinks]{hyperref}
\setlength{\parindent}{0pt}
\setlength{\parskip}{6pt}
\setlength{\emergencystretch}{1em}
\newcommand{\cue}[1]{\par\textbf{#1}\par}
\newcommand{\stage}[1]{\textit{[#1]}}
\raggedbottom
\begin{document}
\begin{center}
{\Large\bfseries Classroom Presentation Script}\\[3pt]
{\normalsize Fourier-Domain Spatial-Mode Conditioning for Coherent Optical Communication}\\[3pt]
Work by Afam Madu and Tofayel Karim
\end{center}
\textit{Read the paragraphs aloud; bracketed directions are actions, not spoken words. Display the paper's figures alongside this script. Figures 1--10 preserve the original baseline; Figures 11--12 present the new investigation. Headings use the titles printed on the plots.}

\cue{Opening story: receiving light is not the same as receiving its message}
Imagine a choir whose singers all know the correct note, but some sing out of step. The sound is still there, yet it does not combine as intended. Light also has a strength and a rhythm. We call its strength pattern \emph{intensity}, and its position within an oscillation cycle \emph{phase}. Our question is whether adjusting this phase can make distorted light fit a receiver better, without adding more light.

The receiver has a reference beam called a \emph{local oscillator}, or LO. Think of it as the conductor's reference rhythm, extended across the beam. A \emph{spatial mode} is the beam's shape and phase pattern. Good reception requires the incoming mode to match the reference. We test this in a computer simulation, not a laboratory experiment.

\cue{Figure 1 --- 4f architecture and model boundary}
\stage{Trace the diagram from source to receiver, then point to the model boundary.}
The source is a smooth, bell-shaped Gaussian beam. Displacement moves it sideways; tilt changes its direction; defocus and astigmatism distort its phase. The first lens maps the beam into a Fourier plane. Like separating a musical chord into notes, this separates a spatial pattern into spatial-frequency components: slowly or rapidly varying features across the beam. These are not different colors of light.

At that plane, a pupil limits what passes, and a phase mask changes the timing of the components. The second lens recombines them. We model this as an ideal Fourier transform and its inverse. The name 4f refers to the usual two-lens relay geometry. The diagram then changes models: optical propagation supplies coupling, and a separate Gaussian-noise model predicts the receiver measurement. We do not simulate every quantum property of the whole apparatus.

\cue{Figure 2 --- Spatial intensity and phase}
\stage{Point left to right across the top row, then across the bottom row.}
The columns show the source, the Fourier plane \emph{before} the mask, and the optimized detector field. The top row shows intensity; the bottom shows phase. The horizontal and vertical axes are positions in millimeters. Read the color bars: intensity color means strength, while phase color means position around an oscillation cycle. Phase is hidden where the field is too weak to interpret meaningfully.

The Fourier-plane and detector-plane scales are different, even though the arrays have equal pixel counts. A dim-looking Fourier panel does not mean energy vanished: power must be summed using each plane's physical pixel area. This is why we checked normalization before trusting any optimization.

\cue{Figure 3 --- Correction comparison and residual phase}
\stage{Compare the three columns: distorted, analytic, optimized.}
The top row shows received intensity. The bottom shows phase relative to the LO after removing one overall phase offset. The analytic mask is a formula-based correction; the optimized mask is learned numerically. Overlap is a shape-matching score between zero and one, separate from how much power survives. For this example, overlap rises from about 0.409 without correction to 0.453 with either correction. Almost all optical power survives these three cases, so the benefit is better matching, not amplification. This is one example, not yet proof that the method works on unseen distortions.

\newpage
\cue{Act II: learning a correction, then giving it an honest exam}
We have shown that correction can help one field. Now we ask how the mask was chosen and whether its benefit survives beyond that example. Think of training as practice problems, validation as a mock exam used to choose a solution, and testing as a final exam that must stay unseen.

\cue{Figure 4 --- Optimized Fourier-plane phase mask}
\stage{Point to the phase map on the left and the histogram on the right.}
The left panel is the learned correction recipe. Position is on the two axes, and color represents phase in radians. Like asking different choir members to shift their timing, the mask changes phase rather than adding amplitude. Phase wraps around a full cycle, so a color jump near minus pi and plus pi need not be a physical discontinuity in the optical action.

The right panel counts how many pixels have each phase value. It includes pixels outside the pupil and pixels receiving little light. A large count therefore does not mean that phase contributes strongly to reception. This is a continuous numerical phase map, not a manufactured device or an etching-depth drawing.

\cue{Figure 5 --- Training and held-out validation}
\stage{Point to iteration on the horizontal axes and overlap on the vertical axes.}
The left panel shows performance on training fields; the right shows performance on separate validation fields. Each curve is a different starting point for the optimizer. Higher overlap is better. There are twelve training distortions, twelve validation distortions, and three optimization starts, each run for one hundred updates. The shaded range shows variation among starts; it is not a statistical confidence interval.

We keep the mask with the best validation score, not simply the last update. The saved best checkpoint is labeled initialization two, iteration fifty-four; those labels count from zero. Its validation overlap is about 0.456. Validation helps choose the mask, so validation performance alone is not an independent final result.

\cue{The final exam: paper Table III, not another learning curve}
\stage{Show Table III or state these results before moving to Figure 6.}
On twenty-four unseen distortions, mean overlap is about 0.431 without correction, 0.470 with the fixed analytic mask, and 0.470 with the trained mask. We compare methods on the \emph{same} distortions, like comparing two runners on the same courses rather than on different hills.

The trained-minus-uncorrected improvement is about 0.039, with a ninety-five-percent confidence interval from 0.0358 to 0.0422. This supports improvement for the tested distribution. Against analytic correction, the interval runs from about minus 0.00137 to plus 0.00191. Because it includes zero, we cannot establish that training is better. That is not proof the methods are identical. Also, the analytic mask stays fixed: it is not recomputed for each test input.

\cue{Figure 6 --- Misalignment and aberration robustness}
\stage{Compare the three heatmaps using their shared zero-to-one color scale.}
The horizontal axis is sideways displacement in micrometers; the vertical axis is defocus strength in radians. Color shows overlap. All other distortion settings stay fixed, and neither mask is retrained at each point. Think of these as terrain maps: they show where a receiver works well and where it struggles, not just its performance at one address.

The result depends on the operating point. These maps are not probabilities, confidence regions, or a guarantee of improvement everywhere. They complement the unseen test examples, but they do not replace them. We have now tested optical matching; next we must check what that matching means for noise.

\newpage
\cue{Act III: the quantum noise floor is not a volume control}
A quieter output is not automatically a better output. Turning a radio down reduces the music as well as its audible loudness. Similarly, loss can reduce our received signal without improving its quality. To avoid this mistake, we first check the noise convention independently of the communication experiment.

Why does this matter for communication? Information is carried by controlled changes in light. If the beam arrives distorted, the receiver can miss part of that information-bearing signal. This simulation lets us compare corrections before building hardware and predict how better matching improves signal-to-noise ratio. It is like adjusting an antenna rather than increasing transmitter power. Such improvements could help an optical link tolerate misalignment or distinguish messages more reliably, but actual error rates, data-rate gains, and secure quantum communication require further testing beyond this study.

\cue{Figure 7 --- Gaussian pure-loss validation}
\stage{Use paper Figure 7, even though its original filename starts with 10. Point to both panels.}
The horizontal axis is transmissivity: zero means none of the input signal passes, and one means all of it passes. The left panel shows the mean of one measured light-wave component, called a \emph{quadrature}. Think of a quadrature as one projection of an oscillation onto a chosen reference rhythm. The right panel shows its variance, meaning the spread of repeated measurements around their mean.

A coherent state is the ideal laser-like reference used here. Loss makes its mean shrink with the square root of transmissivity. But its intrinsic quadrature variance stays at one \emph{shot-noise unit}. That unit is defined by vacuum fluctuations. Even when the mean disappears at zero transmission, the quantum noise floor does not disappear. The choir analogy explains phase alignment, but this vacuum floor is genuinely quantum, not imperfect singing.

The lines show theory, and the crosses summarize one hundred thousand random NumPy samples at each transmission. Small deviations from one are sampling fluctuations, not noise suppression below the quantum limit. Separately, Strawberry Fields, an independent quantum-software backend, checks the means and covariances of twenty-two vacuum and coherent-state cases. Its agreement is near machine precision. Those backend checks are distinct from the plotted random samples. This figure contains pure loss only, not the extra receiver and electronic noise used next.

\cue{Figure 8 --- Calibrated one-quadrature homodyne}
\stage{Point to the five distribution panels, then the bottom-right scatter plot.}
Homodyne detection compares incoming light with the LO to measure one quadrature. Alice is the sender, and Bob is the receiver. Alice varies the prepared quadrature to carry information. Bob sees a reduced version of it plus random noise. Here we assume the overall LO phase is correctly compensated; a real receiver would have to estimate and control that phase.

The five distribution panels represent ideal, distorted, analytic, optimized, and deliberately lossy cases. Their horizontal axis is measured quadrature divided by a separately estimated vacuum standard deviation. Their vertical axis is probability density: taller bars mean values occur more often, not that the beam is brighter. Each case contains one hundred thousand simulated symbols.

Each panel distinguishes three distributions. The unconditional distribution includes all messages and noise. The latent residual isolates the noise because the simulator knows the components it generated; a real receiver does not get that information for free. The conditional distribution selects only symbols for which Alice's value is close to zero. That is like comparing repeated deliveries of nearly the same message, rather than mixing many different messages. The allowed interval has a finite width, so a little signal variation remains. Gaussian curves describe the unconditional distribution and residual, not an exact fit claimed for the finite-bin distribution.

Finally, the scatter plot places Alice's prepared value horizontally and Bob's single measured quadrature vertically. The sloping cloud shows that Bob follows Alice, while its thickness shows noise. It does not mean Bob measured two incompatible quadratures simultaneously. The overall lesson is to distinguish information-bearing variation from unwanted random variation before interpreting a wide or narrow histogram.

\newpage
\cue{Act IV: deciding what improved, and what cannot be recovered}
\cue{Figure 9 --- Variance and uncertainty}
\stage{Point to the case labels, theory crosses, empirical dots, error bars, and vacuum line.}
The horizontal axis lists the five cases. The vertical axis is unconditional quadrature variance in nominal shot-noise units. Unlike the preceding histogram display, these values are not divided by each case's measured calibration estimate. Crosses are theory, dots are simulated estimates, and error bars are approximate ninety-five-percent intervals from finite sampling. Every theoretical value lies within its corresponding interval. This supports consistency of the sampler with the stated model, not agreement with a laboratory measurement.

The ideal case has the largest variance because it preserves the most message variation. The lossy case approaches the noise floor because most useful signal is gone. That smaller variance is \emph{not} better communication or squeezed light. The error bars also omit uncertainty in the vacuum calibration, so they are not a complete experimental error budget.

\cue{Figure 10 --- SNR and efficiency budget}
\stage{Explain the left panel first, then the grouped bars on the right.}
Signal-to-noise ratio, or SNR, compares useful message variation with unwanted noise. On the left, horizontal position is mode overlap and vertical position is theoretical SNR in decibels, a logarithmic scale. Higher is better. In ordinary linear units, the representative SNR rises from about 1.145 without correction to 1.264 with training; analytic correction gives about 1.266. These are representative cases, not the twenty-four-sample test averages.

The right panel explains why overlap alone is not enough. The grouped bars separate optical throughput, mode overlap, channel transmission, detector efficiency, and their product. Think of delivering a parcel: it must survive the journey, reach the correct address, and be accepted. Passing one stage does not compensate for failing another. These efficiencies multiply rather than add, and the plotted cases need not follow one SNR-versus-overlap curve because their throughput differs.

The deliberately lossy case combines half-plane obstruction with attenuation. Only about 6.17 percent of the original power survives the optical system. Even the best unconstrained single-Fourier-plane phase alignment could reach only about 0.244 normalized overlap for this damaged field, with a coupled-power ceiling of about 1.51 percent of the original input \emph{before} channel and detector losses. This ceiling comes from mismatched Fourier magnitudes, not merely from an optimizer that stopped early. Retiming the choir cannot bring back missing singers. Likewise, a phase mask cannot restore lost photons or arbitrarily repair amplitude damage.

\cue{Why trust the calculation, and where does that trust stop?}
\stage{Refer to the paper's verification checklist table.}
We checked Fourier round trips, physical power normalization, a known Gaussian transform, agreement between numerical implementations, and the gradient actually used to train the mask. We checked held-out evaluation, noise statistics, loss bounds, and independent Gaussian-backend moments. Forty-one automated tests passed, and twenty-nine archived artifacts matched their recorded hashes. Hashes check that saved evidence is unchanged; they do not prove the physical model is correct.

Those are the checks for the original archive. They establish numerical consistency, not experimental validation. The new investigation adds eight tests, for forty-nine in total, larger independent tests, and window/pitch convergence checks for analytic architectures and bounds. It does not retrain every mask on every grid. Realistic phase-mask constraints and measured hardware errors remain open, as do any secret-key calculation or security proof.

\cue{A new question, not a replacement for the old evidence}
So far, the figures and numbers came from the preserved baseline. Why did its two corrections end up so close? Was training too short, or were we asking the wrong optical element to do the job? The new investigation keeps the archive intact and tests that question separately. The old twenty-four-input average near 0.470 is not the nominal-field limit near 0.453. They describe different things. We now follow the evidence from the plateau to a different architecture.

\newpage
\cue{Act V: the limitation belongs to one plane, not to all phase optics}
\cue{Figure 11 --- Why optimization plateaus: a single Fourier plane cannot reshape its spectrum}
\stage{Display the new four-panel diagnosis figure; start at panel A.}
The horizontal bars show mean coupled power on the same twelve hundred and eighty fresh inputs. Coupled power, called Q, is the fraction of incident power that reaches the selected target mode: throughput times normalized overlap. It is dimensionless. Black whiskers give ninety-five-percent intervals for each mean, not for paired differences. Gray means no correction; blue is analytic one-plane correction; purple bars are learned corrections; orange is an adaptive bound; teal is the fixed analytic two-plane system we will explain next.

The original analytic mask gives about 0.449 on these fresh draws, rather than its old 0.470 average. The draws changed; the archive did not. Even the orange oracle reaches only about 0.452. An oracle here means an ideal mask recomputed for each input, not a receiver we have built. The fixed two-plane system reaches about 0.949. This large difference points to architecture rather than simply more training.

\stage{Point to panel B, then panel C.}
Panel B asks what the single-plane ceiling becomes when we remove distortion groups from the nominal input, with no random jitter. Horizontal position is overlap ceiling; the rows name the retained distortions. Displacement alone can be corrected, while defocus is the largest isolated limitation. The combinations interact, so these are not separate losses we can add together. In this model the curved phase is centered on the grid origin, not on the shifted beam. A displaced beam samples a slope in that curved phase, adding spectral shift. Centering the curvature on the beam would be a different physical model, not a correction to the archived data.

Panel C makes the mechanism visible. Horizontal position is the Fourier-plane coordinate in millimeters. Height is normalized intensity integrated along the other axis, in inverse millimeters. The red distorted spectrum is shifted and broadened compared with the blue target. A mask placed right there can retime those spatial-frequency components but cannot change their strengths. It is like conducting a chord whose notes have the wrong relative loudness: timing alone is insufficient. For the nominal field, analytic phase alignment already reaches the exact one-plane ceiling, about 0.453. No extra training in that same plane can exceed it.

The dashed teal curve is different: phase was corrected at the input \emph{before} the transform. Its spectrum matches the target. Changing phase in one plane can change intensity after propagation to another. That is why two phase screens can succeed without either one amplifying light or violating the original bound.

\stage{Point to panel D and read its legend.}
The horizontal axis multiplies all jitter standard deviations; vertical position is mean coupled power. The blue fixed-one-plane and teal fixed-two-plane curves are compared with the orange adaptive one-plane bound and purple empirical fixed-unitary ceiling. These use a separate two-hundred-and-fifty-six-input diagnostic set, not the primary test. A unitary system is an ideal lossless transformation that preserves the differences between input modes. It cannot turn every different random input into exactly the same output. The empirical ceiling quantifies that restriction for the finite sampled ensemble; it is not a population confidence band. Jitter, rather than the old single-plane bottleneck, now limits the fixed two-plane receiver.

\newpage
\cue{Act VI: put each correction where it can do the job}
\cue{Figure 12 --- Architecture change: correct phase before the Fourier plane, then recenter the beam}
\stage{Trace the top schematic from left to right, including its final relay note.}
First place a phase screen at the input, or at an optical image of that input. It cancels the mean tilt, defocus, and astigmatism. The first lens then makes the Fourier transform. At its Fourier plane, a phase ramp removes the mean sideways displacement. Finally, the inverse relay delivers the beam to the same target mode and LO. The first correction prevents spectral shift and broadening; the second recenters the beam. Both masks stay fixed during testing. They are calculated analytically, not learned and not adjusted separately for each test input.

\stage{Point to panel A, then panel B.}
Panel A places one-plane coupled power horizontally and two-plane coupled power vertically. Each dot compares the same distortion, and the dashed diagonal means equal performance. On the fresh test, mean Q rises from 0.449020 to 0.948633. The paired gain is about 0.4996, with a ninety-five-percent interval from 0.4950 to 0.5042. This is strong evidence for the architecture change within the assumed distribution, not a laboratory demonstration. Estimating the mean correction from training inputs instead of using the known nominal parameters still gives about 0.9473.

Panel B asks whether adding hardware would waste the gain. Its horizontal axis is extra power transmission, and its vertical axis is mean paired gain after that extra loss. The shaded band is an approximate pointwise confidence interval. Mean break-even occurs near 0.473 transmission, equivalent to about 3.25 decibels of additional loss. At ninety percent extra transmission, mean Q is about 0.854. These are sensitivity calculations, not measured component efficiencies, and mean break-even does not guarantee every individual input improves.

\stage{Point to panel C and emphasize its small axis scale.}
Panel C revisits training with matched controls. Its horizontal units are ten to the minus four in coupled-power gain. With twelve training inputs, validation retains the analytic mask. With one hundred and twenty-eight, it finds a tiny gain, about 0.000125. One hundred and eight hundred updates choose the same checkpoint. Changing to a power or conditional-information objective does not unlock a large gain. These controls were declared after the primary run and are exploratory. There are no error bars here; the report contains their intervals. The original trained mask is actually slightly worse than analytic on the fresh test. Thus we should not say training can never help: we should say the gains found here are tiny beside the architecture change.

\stage{Point to panel D, then return to the schematic.}
Panel D compares five independent seed blocks, each with two hundred and fifty-six inputs. Vertical position is mean coupled power or its empirical bound. The two-plane result stays close to each block's fixed-unitary ceiling. Those five bounds are not a confidence interval. The nominal field is corrected almost perfectly, but randomly varying fields cannot all be corrected perfectly by one fixed receiver. Further large gains would need stabilization, adaptive estimation, or detection of multiple spatial modes, with a new measurement and noise budget.

\cue{Closing statement: a physical diagnosis, not a claim that learning wins}
The original baseline showed that phase correction helps, without demonstrating trained superiority. The new evidence explains the plateau: a single Fourier-plane mask cannot reshape the spectral magnitudes it receives. Correcting phase upstream, then displacement in the Fourier plane, changes that limit and roughly doubles useful coupling in this simulation. The winner is an analytic two-plane design. It neither restores lost photons nor lowers vacuum noise. The original homodyne plots remain baseline plots; they are not new two-plane measurements. The useful next experiment is to calibrate both phase planes and measure coupled power and transmission on the same perturbations, before claiming hardware, data-rate, or quantum-security gains.
\end{document}
